\documentclass[11pt]{article}
\usepackage[margin=1in]{geometry}
\usepackage{amsmath,amssymb,amsthm}
\newtheorem{theorem}{Theorem}
\newtheorem{lemma}{Lemma}
\title{Analytic Existence of the First Postcritical Contact Branch}
\author{}\date{October 5, 2026 --- paper proof}
\begin{document}\maketitle
Use the explicit first-ordering formulas of \texttt{contact-equations.tex}, with $w=1+\delta$ and $0<\phi<b<c<w/2$. Let $R_1,R_2,R_3$ be the three shooting residuals in that note, in the order printed there.

\begin{theorem}[A unique local analytic branch]
There is $\varepsilon>0$ and a unique analytic branch of solutions in a neighborhood of the following scaled parameters:
\[
 \phi=\delta^2\lambda(\delta),\qquad b=\delta\mu(\delta),\qquad c=\delta\nu(\delta),
\]
\[
 \lambda(0)=3/16,\qquad\mu(0)=3/2,\qquad\nu(0)=2.
\]
For $0<\delta<\varepsilon$ the ordering $0<\phi<b<c<w/2$ holds. The corresponding supports define a normalized cap with a directly feasible cap-minus-niche sofa. This theorem concerns the shooting branch and feasibility; its exact niche boundary and global optimality are proved separately.
\end{theorem}
\begin{proof}
The elementary trigonometric and integral expressions in \texttt{contact-equations.tex} extend analytically as functions of $(\delta,\lambda,\mu,\nu)$ near $(0,3/16,3/2,2)$. This refers to their explicit formulas with oriented integrals, not to interpreting an invalid negative-$\delta$ contact ordering as a geometric construction. The two denominators determining $g_0$ and $p_0$ tend respectively to one and $\cos(1/2)+\sin(1/2)>0$.

The desingularized residuals $(\delta^{-3}R_1,\delta^{-2}R_2,\delta^{-2}R_3)$ extend analytically to $\delta=0$. Their values there are
\begin{align*}
 r_1&=2\lambda\mu-\mu^2/2+\mu^3/6,\\
 r_2&=2\lambda-\mu+\mu^2/2,\\
 r_3&=\nu-\nu^2/2-\lambda-d(\mu,\nu),\qquad
 d(\mu,\nu)=\frac{(\nu-\mu)(2-\mu-\nu)}4.
\end{align*}
Here is a direct way to check the cancellation, including the initially unknown $p_0$. On the first interval,
\[
 \gamma_x(\phi)=p_0-g_0\sin^2\phi-\sin\phi\cos\phi-\cos\phi+\sin\phi,
\]
\[
 \gamma_y(\phi)=g_0\sin\phi\cos\phi+\cos^2\phi-\sin\phi-\cos\phi.
\]
The $p_0$ terms cancel from $R_1,R_2$. The remaining support integrals at $b$ are those of $A-1-t$ on $[\phi,b]$. At $c$, the floor residual is exactly
\[
 -\sin c+\int_\phi^b(A-1-t)\cos t\,dt
             +\frac12\int_b^c(A-t)\cos t\,dt.
\]
Using $g_0=2+\delta+O(\delta^2)$, $A=2+\delta+O(\delta^2)$, $\phi=\lambda\delta^2$, $b=\mu\delta$, $c=\nu\delta$ gives the three polynomials above. All lower Taylor coefficients vanish identically in the scaled variables, which justifies the analytic division by the displayed powers of $\delta$.

At $(\lambda,\mu,\nu)=(3/16,3/2,2)$ the three polynomials vanish and their Jacobian is
\[
 \begin{pmatrix}3&0&0\\2&1/2&0\\-1&-1/4&-1/2\end{pmatrix},
 \qquad \det=-3/4\ne0.
\]
The analytic implicit function theorem gives the stated unique local branch. Positivity and the required ordering follow from its limiting scaled parameters.

We verify admissibility instead of treating the equations as sufficient by themselves. The initial arms satisfy $f(0)=1$ and $g(0)=2+O(\delta)$. Before $\phi$ they satisfy $f'=g$, $g'=-f$; on the following core-only intervals they satisfy $f'=1$, $g'=-1$; on the right-envelope interval they satisfy $f'=g/2$, $g'=-1$. Reflection supplies the second half. Hence $f>1$, $g>1$ in the interior, $f'>0$, and $g'<0$ for sufficiently small $\delta$.

The cap curvature densities from the balance equations are nonnegative. The supports and their first derivatives converge to those of the angle-one relaxed cap. On a fixed sufficiently small interval near zero one has $p''<-(\sec1-\tan1)/2$ in every piece, while $p'(0)=0$. Away from that interval the limiting support has strictly negative derivative. Consequently $p'<0$ on $(0,w)$ and, by reflection, $k'>0$. Since $p(w)=k(0)=1$, all interior thresholds are positive.

Extend the active supports over the complementary normal intervals by the support of the fixed upper vertex and by the cosine pieces determined by $p(0)$ and $k(w)$. The derivatives at the outer active endpoints are zero. The upper pinned edge lengths tend to one and the lower fan-edge lengths are positive. Thus the full circular support is continuous and has nonnegative curvature measure, so defines a normalized cap. This step includes the necessary normal-gap extensions when the unpinned endpoint values change.

Finally, its fan-truncated niche quadrilaterals converge uniformly to those at the angle-one relaxed cap, whose vertices lie in a disk of radius less than one half. They therefore lie in a common disk of radius $\rho<1$. The cap contains the unit fan sector, its niche is star-shaped and relatively open in the fan, and radial paths connect its complement to the unit fan arc. The canonical supporting-hallway motion is feasible. This proves compactness, path-connectedness and feasibility of the cap-minus-niche sofa, without yet identifying every exposed boundary arc.
\end{proof}

\begin{lemma}[The floor extends past the initial corner]
Along this branch,
\[
 \frac{p(c)-1}{\cos c}-(p(0)-1)=\frac13\delta^3+O(\delta^4).
\]
\end{lemma}
\begin{proof}
The $p_0$ term cancels from the left side, leaving
\[
 1-\sec c+\frac1{\cos c}\left[
 \int_\phi^b(A-1-t)\sin(c-t)\,dt+
 \frac12\int_b^c(A-t)\sin(c-t)\,dt\right].
\]
Its cubic coefficient is
\[
 \nu^2/2-\nu^3/6-\lambda\nu
 +\frac12\int_\mu^\nu(s-1)(\nu-s)\,ds.
\]
At $(\lambda,\mu,\nu)=(3/16,3/2,2)$ this is $1/3$. Analyticity supplies the stated remainder.
\end{proof}

\paragraph{Scope and validation.}
The positive interval supplied here is existential; no numerical lower bound for $\varepsilon$ is asserted. The Jacobian is an exact rational certificate, not a floating-point nonsingularity test. Root existence near the transition is not being extrapolated across the later changes in contact order. No CI or compilation was run.
\end{document}
